\chapter{Method and Implementation}\label{ch:method}

This chapter describes \emph{how} the study was built: the from-scratch State-Optimization
predictive-coding (SO-PC) network and the single interface through which every experiment runs;
the correctness regime that lets us trust hand-derived updates; the fused CUDA settling kernel
that is the project's principal systems contribution; the fair-comparison protocol that
disciplines the predictive-coding-versus-backpropagation (PC-vs-BP) question; the
continual-learning harnesses (including a Song-exact alternating replication); the generative
variant of the same machinery; the optional incremental-PC and elastic-weight-consolidation
extensions; and the reproducibility tooling. Throughout, the emphasis is on design decisions and
their justification. Quantitative outcomes are deferred to \Cref{ch:experiments}; the underlying
PC mathematics is developed in \Cref{sec:so-pc} and is referenced here rather than re-derived.

A single conviction shapes every choice below. The point of the project is to \emph{understand and
demonstrate} a learning rule that is local --- one in which every parameter and state update can be
computed from quantities available at the two layers it connects, with no global error signal
propagated backward through the network. We therefore use PyTorch \citep{paszke2019pytorch} purely
as a tensor engine: there is no call to \texttt{loss.backward()} anywhere on the PC path, and every
state and weight update is performed by hand under \texttt{@torch.no\_grad()}. Autograd appears in
exactly two places, both deliberate: as an external oracle in the gradient-correctness tests
(\Cref{sec:method-correctness}), and as the learning rule of the backpropagation baseline against
which PC is measured. This is not an aesthetic preference. If autograd silently supplied a gradient
anywhere in the settling or learning loop, the claim ``this is local learning'' would be hollow, and
the comparison with backprop would be circular. Keeping the PC arm autograd-free is what makes the
study honest.

\section{From-Scratch PCN Architecture and Interface}

\subsection{The model object}

The network is a small object, \texttt{PCN}, holding a list of weight matrices $W_i$, bias vectors
$b_i$, an activation function $\phi$ together with its derivative $\phi'$, and bookkeeping for device
and dtype. For a topology with layer sizes $[n_0, n_1, \dots, n_L]$ (input $n_0$, output $n_L$) there
are $L$ weight matrices, where we write $L = n$ for the number of layers/transformations in the code.
The single primitive the rest of the system relies on is the per-layer prediction
\[
  \mathrm{pred}_{i+1} = W_i\,\phi(s_i) + b_i,
\]
implemented as \texttt{predict(i, state)}. Everything else --- feedforward initialisation, settling,
the local weight update, evaluation, the kernel, the baselines --- is expressed in terms of this one
function, so that the forward computation is defined in exactly one place and cannot drift between
the PC arm and its backprop control.

Weights are initialised orthogonally by default. This is a substantive choice rather than a default
inherited from convention: orthogonal initialisation keeps the layer-to-layer Jacobian eigenvalues
near one, which mitigates the exponential signal decay that afflicts deep SO-PC during settling
(\Cref{sec:so-pc} and the depth discussion below). The activation is \texttt{tanh} by default, with
\texttt{identity} and \texttt{sigmoid} also supported; \texttt{sigmoid} matters because it is the
activation Song et al.\ use in the regime we replicate (\Cref{sec:method-continual}).

\subsection{Feedforward initialisation, settling, and the local update}

Training proceeds per minibatch in four conceptual stages, matching the SO algorithm
of \Cref{sec:so-pc}:

\begin{enumerate}
  \item \textbf{Feedforward init.} The input is clamped at $s_0 = x$ and the remaining states are
  filled by a single forward pass, $s_{i+1} \gets \mathrm{pred}_{i+1}(s_i)$. By construction this
  sets every internal error $\varepsilon_k = s_k - \mathrm{pred}_k$ to zero, so settling begins from a
  zero-internal-energy state and only the clamped output injects a non-zero error. This is not merely
  a convenient warm start: it coincides with the first step of the error-optimization view of PC, which
  is part of why it is such a good initialiser.
  \item \textbf{Clamp.} The input node is clamped to the image and (during training) the output node
  is clamped to the one-hot label.
  \item \textbf{Inference / settling.} With the weights held fixed, the free states are relaxed for up
  to $T$ steps toward the energy minimum. Each step is a synchronous (Jacobi) update: all errors
  $\varepsilon_k$ are computed first, then all free-state updates
  $s_k \gets s_k - \eta_x\big(\varepsilon_k - \phi'(s_k)\odot (W_k^{\!\top}\varepsilon_{k+1})\big)$ are
  applied from those errors, so that no layer's update sees another layer's change within the same step.
  The output layer, when free at test time, updates with $\varepsilon_k$ alone.
  \item \textbf{Learning.} At the settled equilibrium the weights are updated by the purely local
  Hebbian rule $W_i \mathrel{+}= \eta_w\,\varepsilon_{i+1}\,\phi(s_i)^{\!\top}/B$ and
  $b_i \mathrel{+}= \eta_w\,\overline{\varepsilon_{i+1}}$, with $B$ the batch size. Each update uses only the
  post-synaptic error and the pre-synaptic activation of the layer it belongs to.
\end{enumerate}

The settling routine is deliberately compact --- on the order of fifty lines --- because it is the
conceptual core of the repository and must be auditable line by line. It takes flags that, in the same
few lines, support the discriminative training case (input and output clamped, hidden states free) and
the generative / inpainting / anomaly cases (input free or partially masked), which is what lets the
generative experiments reuse exactly the same dynamics (\Cref{sec:method-generative}).

\subsection{The convergence criterion}

By default settling runs a fixed number of steps $T$. We added an optional tolerance: when a relative
energy-change threshold \texttt{tol} is supplied, the loop stops once
$|E_t - E_{t-1}| / (|E_{t-1}| + \epsilon) < \texttt{tol}$. With \texttt{tol = None} the original
fixed-$T$ behaviour is reproduced exactly. This small addition turns ``how many steps does settling
need to converge?'' into a measurable quantity rather than a hyperparameter we assume. We use it as a
diagnostic: \texttt{settling\_steps\_to\_converge} is measured on a training batch with a generous step
cap (five times $T$, at least one hundred) and a reference tolerance of $10^{-3}$, so the reported value
is the true number of steps the dynamics need to plateau, independent of the configured $T$. If that
number exceeds the training $T$, the network is under-settling --- directly relevant because PC's
claimed advantages are predicated on reaching the activity equilibrium \citep{song2024prospective,
innocenti2026limits, qi2025settling}.

\subsection{The \texttt{train\_and\_eval(config)} contract}

The only deliberate abstraction in the system is a single function, \texttt{train\_and\_eval(config)
$\rightarrow$ metrics}. It accepts a configuration dictionary --- hidden sizes, $T$, the state and
weight learning rates, activation, weight init, number of epochs, batch size, seed, backend, and so on
--- builds a PCN, trains it on MNIST \citep{lecun1998mnist}, and returns a metrics dictionary
containing test accuracy, an energy curve (mean equilibrium energy per epoch), noise-robustness
(accuracy across input-noise levels), training wall-clock, the settling-steps-to-converge diagnostic,
and the fully resolved configuration. There is no plugin or adapter architecture: a single coherent
implementation with one mechanism, exposed through one interface. The autonomous hyperparameter search
(\Cref{sec:method-repro}) and the experiment harnesses are thin layers that call this function and log
the result --- no rewrite, by design.

Two small details in the contract guard against silent error, in keeping with the honesty principle.
First, the thesis-specification names $\eta_x$ and $\eta_w$ are accepted as configuration aliases for
the internal \texttt{lr\_state} and \texttt{lr\_weight}, with the spec names taking priority, so a config
written to the documented notation is honoured rather than silently ignored. Second, unknown keys raise
a warning instead of being dropped, and reserved-but-unimplemented features (non-isotropic precision
schedules; update variants other than the implemented two) raise \texttt{NotImplementedError} rather
than silently no-op'ing. The default weight learning rate was raised from $10^{-3}$ to $10^{-2}$ after
the lower value was found to underfit MNIST; precision is hard-coded to the isotropic case $\Pi = I$,
which reduces the energy $E = \tfrac12\sum_k \lVert\varepsilon_k\rVert^2$ to a plain sum of squared
errors and is mathematically exact for that choice.

\section{Correctness: Hand-Derived Gradients vs.\ Autograd}\label{sec:method-correctness}

Because no autograd runs on the PC path, the burden of proof falls entirely on us: the settling update
and the Hebbian weight update are gradients of the PC energy that we derived and transcribed by hand,
including the easily-mistaken signs, the transpose on the feedback term $W_k^{\!\top}\varepsilon_{k+1}$,
and the $\phi'(s_k)$ gating factor. A sign error or a missing transpose would not crash --- it would
quietly train a slightly different objective and contaminate every comparison. We therefore treat
correctness as a first-class deliverable and pin it down two ways.

First, against an analytic reference. The exact SO update --- signs, transposition, and the $\phi'$
gating --- is validated against an independent analytic NumPy derivation, and the implementation is
held to match it; this is a hard convention of the project, not a one-off check. Second, against
autograd as an external oracle. In the test suite (and only there) we build the same energy with
autograd enabled and confirm that the hand-coded settling-state gradient and the hand-coded weight
update agree with \texttt{torch.autograd} to roughly $10^{-5}$. This is the one legitimate use of
autograd on the PC side: not to \emph{compute} the updates the network learns by, but to \emph{certify}
that our local rules equal the true energy gradients. A complementary property test asserts that the PC
energy decreases monotonically under settling on toy tensors, so that a regression in the dynamics is
caught even when the analytic check passes. The tests are hermetic --- they run on small synthetic
tensors with no MNIST download and complete on CPU --- so that correctness is cheap to re-verify on any
machine (\Cref{sec:method-repro}). The exact derivations these checks target are given in
\Cref{app:derivation}.

\section{The Fused CUDA Settling Kernel}\label{sec:kernel}

\subsection{Motivation: predictive-coding inference is launch-overhead-bound}

The defining cost of PC, relative to a feedforward network, is the iterative inference phase: each
training (or test) example requires $T$ settling steps before any weight update. On hardware optimised
for the single dense matmul pattern of backprop, this iterative structure is a mismatch. Crucially,
within a single settling step the work is not one large operation but a stream of \emph{small} ones ---
a small matmul to form each prediction, an elementwise subtraction for each error, a transpose-matmul
for the feedback term $W_k^{\!\top}\varepsilon_{k+1}$, the elementwise $\phi'$ gating, and the state
update --- and these recur for every layer and every step. At MNIST widths (input $784$, hidden $256$)
each of these operations is tiny, so the dominant cost is not arithmetic but the fixed overhead of
launching a GPU kernel for each of them. A standard PyTorch settling step issues many such launches; we
measure exactly $31$ kernel launches per step on our two-hidden-layer topology, independent of width,
so that $T=20$ costs $620$ launches and $T=80$ costs $2480$. This is the ``framework tax'': the wall
clock is set by launch count, not by floating-point work. It also explains the low GPU utilisation
(roughly a quarter to a half) we observe for the PyTorch PC arm.

The opportunity is precisely the parallel structure of one settling step. Within step $t$, all errors
$\varepsilon_k$ depend only on the states at $t$ and can be computed together; then all free-state
updates $\Delta s_k$ depend only on those errors and can be applied together. Only the chain of $T$
steps, and the feedforward initialisation, are inherently sequential. A handwritten kernel can therefore
fuse an entire step --- or the whole $T$-step loop for a batch --- into a single launch, holding the
states resident in shared memory or registers across all $T$ steps and reading the weights from global
memory once (persistent-kernel style). The number of launches collapses from hundreds or thousands to
\emph{one}.

\subsection{Novelty and honest scope}

We audited the open-source landscape to place this contribution honestly. Across eleven open PC
libraries there is no handwritten CUDA inference kernel: they are uniformly JAX/JIT-based or built on
PyTorch autograd, and contain zero \texttt{.cu} files. We extended the audit to the closely related
Equilibrium Propagation family (eighteen repositories) and found the same gap, the nearest neighbour
being a per-step Triton kernel driven by a host-side loop. The defensible claim is therefore narrow and
specific: \emph{the first handwritten fused CUDA settling kernel for predictive coding, to our
knowledge}. We are equally explicit about what is \emph{not} novel. The fusion technique itself ---
persistent, state-resident, single-launch kernels --- is established prior art from persistent RNNs and
large-language-model megakernels, and the ``launch-overhead-dominated'' observation is a generic
framework property, not our discovery. Our contribution is the transfer of that technique to PC plus a
clean measurement of how much of PC's inference overhead is launch overhead. The most recent
algorithmic neighbour, error-optimization PC \citep{goemaere2025epc}, addresses settling efficiency at
the algorithm level rather than the kernel level and is cited as the contrasting approach. The kernel
audit and verdict are documented in full as part of the project record.

\subsection{The fused design and its three versions}

The CUDA source is the single source of truth: a real \texttt{.cu} file compiled at first use via
\texttt{torch.utils.cpp\_extension.load} (JIT, cached), not an inline string or a commented stub. The
C++/CUDA extension exposes the kernel through a \texttt{pybind11} module in the same file. Three versions
of the two-hidden-layer kernel evolved, each addressing a specific bottleneck revealed by measurement:

\begin{description}
  \item[v1 --- one block per sample.] Each thread block settles one sample, holding that sample's
  states, errors, and activations in shared memory across all $T$ steps. This eliminates the launch
  overhead and wins decisively at small batch. Its weakness is that every block re-reads the full
  weight matrices from global memory each step, so at large batch it is bandwidth-bound: with one block
  per sample the same weight element is fetched once per sample.
  \item[v2 --- one block per tile of samples.] To recover the lost weight reuse --- the key optimisation
  of any GEMM --- v2 assigns a tile of \texttt{TB}$=8$ samples to each block and reads each weight
  element once, reusing it across the tile via an inner unrolled loop over the tile. The batch is
  zero-padded to a multiple of the tile size, and $\phi(\text{input})$ is precomputed on the host. This
  pushes the crossover with cuBLAS toward larger batches.
  \item[v3 --- cached input activation.] Because the input is clamped, $\phi(\text{input})$ is constant
  across all $T$ steps, yet v1/v2 re-read it each step. v3 caches the $\phi(\text{input})$ tile once in
  shared memory (opting in to the larger per-block shared-memory budget on the target architecture),
  eliminating the roughly $256$-fold redundant global reads of the input layer. This lifts the
  mid-batch speed further and moves the break-even point with cuBLAS to around batch $1024$--$1500$.
\end{description}

The runtime dispatches by batch size: small batches take the per-sample path (v1), larger batches the
tiled path (v2/v3). The activation is selected by an integer code ($0$ tanh, $1$ identity, $2$ sigmoid)
matched between the Python wrapper and the device functions; the wrapper detects the activation from
$\phi(0)$ and $\phi(1)$ and \emph{raises} on anything unrecognised. This guard was added after a real
bug in which sigmoid silently fell through to tanh because the detector keyed only on tanh's signature
--- precisely the kind of silent-wrong-answer failure the project is built to avoid.

\subsection{The general-depth variant}

The v1/v2/v3 kernels hardcode the two-hidden-layer topology for speed, which was a documented
limitation. We removed it with a general-depth kernel that settles a PCN of \emph{arbitrary} depth $L$.
The mechanism is to pass the weights, states, and biases not as fixed arguments but as arrays of device
pointers (carried as 64-bit integer addresses), together with per-layer sizes and \emph{dynamically
computed} shared-memory offsets, so a single kernel can index any number of layers laid out contiguously
in shared memory. The update is the identical Jacobi step used by the PyTorch reference. The dispatch
keeps the optimised two-hidden-layer path for that common case and routes any other depth to the
general kernel. It is per-sample only --- the launch-bound small-batch regime the kernel targets --- so
a register-blocked tiled GEMM for large batch at large depth remains open. A platform pitfall worth
recording: on Windows \texttt{long} is $32$-bit, so the pointer arrays must be explicitly $64$-bit
integers; using the wrong type produces a size/type mismatch. The kernel is correctness-verified across
depths two through five (one to four hidden layers) crossed with tanh and sigmoid and two batch sizes
--- sixteen configurations, all matching the PyTorch path to roughly $10^{-6}$.

\subsection{Correctness before speed}

Following good benchmarking discipline, correctness was established before any timing: the kernel must
reproduce the PyTorch reference, and only then is it timed. Verification uses \texttt{allclose} at
roughly $10^{-6}$ against \texttt{\_settle\_pytorch} across activations (tanh, identity, sigmoid), clamp
on and off, and both dispatch paths. One subtlety is load-bearing: the comparison must be run in a
\emph{stable} settling regime. An unstable configuration amplifies ordinary floating-point reordering
differences between the two implementations into apparent ``errors'' that are really divergence of the
dynamics, not a kernel bug. All timing measurements call \texttt{torch.cuda.synchronize()} before
reading the clock, since CUDA execution is asynchronous and naive timing would measure launch latency
rather than completion. The kernel is strictly optional: the default backend is PyTorch, fully
functional, and the CUDA path is enabled only deliberately, by setting an environment flag and
requesting \texttt{backend="cuda"}; without both, requesting the CUDA backend raises rather than
silently falling back. The kernel speedups, the direct launch-count evidence, and the end-to-end
integration are reported in \Cref{sec:exp-kernel}; the launch-count and speedup figures are presented
there.

\section{The Fair-Comparison Protocol}\label{sec:protocol}

The PC-vs-BP question is easy to answer badly. A casual comparison varies the learning rule \emph{and}
a dozen incidental factors at once, and then attributes the difference to the learning rule. Our
protocol exists to isolate the one variable of interest, following the fairness discipline of
\citet{song2024prospective}.

\subsection{Identical architecture, initialisation, and forward function}

The backprop baseline, \texttt{BPMLPRef}, is constructed by \emph{building a PCN at the same seed and
cloning its weights and biases} into trainable parameters. This guarantees, by construction rather than
by careful matching, that PC and BP start from bit-identical initial weights. The baseline's forward
pass reuses the \emph{same} layer chain as the PCN's feedforward initialisation
($h \gets \phi(h)\,W_i^{\!\top} + b_i$), so the two arms compute the identical function of their
parameters; only the learning rule differs. Both arms draw from the same MNIST pipeline. The backprop
arm \emph{does} use autograd --- that is the entire point, it is the backprop learner --- and returns
the same metrics dictionary schema as \texttt{train\_and\_eval}, so the search loop and comparison
harness treat the two arms uniformly.

\subsection{Matched loss --- the dominant confound}

The single most important control is the loss function. PC clamps the output to a one-hot target, which
makes its implicit objective a squared error to that one-hot vector. A backprop baseline trained with
cross-entropy is therefore \emph{not} a like-for-like control: it differs from PC in both the learning
rule and the objective. The protocol provides both BP arms: BP(MSE), which minimises the same squared
error to the one-hot target that PC implicitly does and is the fair comparison, and BP(CE), the
practical cross-entropy baseline reported for reference. The difference BP(CE) $-$ BP(MSE) then isolates
the pure loss effect. This matters because, as \Cref{sec:exp-pcbp} reports, the loss function turns out
to move the comparison far more than the learning rule does; an early, large apparent BP advantage was
almost entirely the loss, not the rule.

\subsection{Joint learning-rate tuning and the stability frontier}

A second confound is asymmetric tuning. Sweeping only the weight learning rate $\eta_w$ while fixing the
state learning rate $\eta_x$ systematically penalises PC, because the two interact through a stability
frontier: at a given $\eta_x$, increasing $\eta_w$ past a point causes the settling dynamics to diverge,
collapsing accuracy to chance; lowering $\eta_x$ restores stability and lets the same $\eta_w$ succeed.
The protocol therefore tunes $\eta_x$ and $\eta_w$ \emph{jointly} for PC over a grid, and tunes each
method's learning rate independently. A related guard, learned through a false-positive, is that the
backprop baseline must itself actually be \emph{learning} before any comparison is meaningful: an
untuned BP(MSE) that sits at chance can manufacture a spurious PC ``win'' that a learn-accuracy
diagnostic immediately exposes.

\subsection{Statistics: validation selection and bootstrap intervals}

Learning rates are selected on a held-out validation split and reported on the test set, never selected
on the test set. Every headline number is reported as a mean over at least three seeds with a bootstrap
confidence interval, so that ``A beats B'' is a statement about non-overlapping intervals rather than
single-run noise. We are explicit about the interval width used: following Song, the bootstrap routine
reports $68\%$ (one-sigma) intervals, and we label them as such rather than implying $95\%$. The harness
that implements bootstrap intervals and grid sweeps lives in the experiment protocol module and feeds
the comparisons of \Cref{sec:exp-pcbp}.

\section{Continual-Learning Setups}\label{sec:method-continual}

Continual learning is where PC's prospective-configuration advantage is most often claimed, so it
receives the most careful protocol design. All setups share the fair-comparison machinery above ---
identical architecture and initialisation, matched loss, per-method tuning --- and add the canonical
continual-learning metrics. We implement these in-house rather than via Avalanche
\citep{lomonaco2021avalanche} to avoid dependency conflicts and keep the project genuinely
from-scratch.

\subsection{Metrics: the R-matrix, ACC, BWT, and learn-accuracy}

After training task $i$ we evaluate on every task $j$, filling the accuracy matrix $R[i][j]$ in the
sense of \citet{lopezpaz2017gem}. From it we read the average final accuracy
$\mathrm{ACC} = \mathrm{mean}_j R[T-1][j]$ and the backward transfer
$\mathrm{BWT} = \mathrm{mean}_{j<T-1}\big(R[T-1][j] - R[j][j]\big)$, where a negative BWT is forgetting.
Critically, we also report \emph{learn-accuracy}, the mean of the diagonal $R[j][j]$ --- the accuracy on
each task right after it is learned. This third metric exists to defeat the stability-plasticity
confound: ``PC forgets less'' is only meaningful if PC also \emph{learned} each task comparably well. A
method that learns each task poorly will trivially appear to forget little, and learn-accuracy makes
that visible.

\subsection{Domain-incremental: Permuted-MNIST}

Permuted-MNIST is domain-incremental: a fixed ten-way head throughout, with each task applying a fixed
random pixel permutation. Because the head and architecture never change, PC and BP run through an
identical task stream with no head surgery. Task zero is the identity permutation (plain MNIST). The PC
arm trains each task with the standard settle-then-update epoch; the BP arm with SGD on its chosen loss;
both are evaluated on every task after each, capping the per-task test set because a PC evaluation is
itself a settle and the full quadratic evaluation grid would be costly.

\subsection{Class-incremental: Split-MNIST and Split-FashionMNIST}

The class-incremental setups partition the labels into disjoint class sets that \emph{share} a single
$K$-output head, with labels remapped to $0\ldots K-1$ so the shared head is forced to overwrite ---
this interference is exactly the regime in which a PC advantage is claimed, unlike the milder
domain-incremental case. Split-FashionMNIST (two disjoint five-class tasks) reproduces the shape of
Song's setup; Split-MNIST (five disjoint two-class tasks) is the other standard variant. Both
\citep{xiao2017fashionmnist, lecun1998mnist} run through the same harness.

\subsection{The Song-exact alternating protocol}

The setups above train each task to completion before moving on. Song's Figure 4 protocol is different
and more demanding: two disjoint five-class FashionMNIST tasks share one five-output head and are
trained by \emph{minibatch-level alternation} --- a few minibatch updates on the current task, then a
switch, repeated. We implemented this faithfully: a network of four layers with thirty-two hidden units
per layer, sigmoid activation, Xavier-normal initialisation, batch size thirty-two, switching every four
updates. On the axes that define the protocol --- four layers, thirty-two hidden, sigmoid, Xavier, two
disjoint five-class tasks, shared head, alternation --- this matches Song exactly, in contrast to our
earlier sequential proxy.

Two methodological points are essential to report honestly here, both surfaced by an adversarial
three-lens review (literature fidelity, statistics, and alternative explanations) that corrected a
premature reading of the protocol. First, a framing correction: Song's Figure 4e is \emph{not} a
``PC is robust across learning rates'' claim --- that is a different figure. Figure 4e shows \emph{less
catastrophic interference}, plotting the mean test error of both tasks against learning rate \emph{with
the learning rate tuned separately per method}. We therefore test interference, not learning-rate
robustness, and we tune the learning rate per method. Second, a faithfulness-versus-trainability
tension: Song's exact budget of roughly eighty-four minibatch iterations leaves our network at chance.
Rather than silently enlarge the budget, we make the training budget an explicit, controlled axis,
sweeping from the exact eighty-four up to convergence. The reporting metrics distinguish the mean of
both tasks' final accuracy (Song's Figure 4e metric) from the \emph{minimum} of the two --- the worse
task being the true interference indicator, since a high mean with a low minimum is a sacrificed task,
not balanced learning --- alongside learn-accuracy and retention. This protocol is also where the CUDA
kernel earns its keep as a real engine: the small batch size of thirty-two is exactly the launch-bound
regime the kernel targets, and the alternating study can be run on \texttt{backend="cuda"} with per-seed
identical results (\Cref{sec:exp-song}).

\section{Generative PC}\label{sec:method-generative}

A distinctive property of predictive coding is that one mechanism --- settling --- can classify,
generate, inpaint, and detect anomalies, depending only on what is clamped. We implement this without a
new learning algorithm, by reusing the same PCN, settling, and weight update.

\subsection{Orientation flip and generative training}

The discriminative network maps image to label, and is therefore not a generative model of images:
settling ``backward'' from a clamped label finds no learned image manifold and produces noise, as we
verify honestly in \Cref{sec:exp-gen}. The fix is to flip the orientation. A generative PCN maps
\emph{label to image}, with the topology $[10, h, h, 784]$. It is trained generatively --- the label is
clamped at the input node, the image at the output node, the hidden states settle, and the same local
Hebbian update is applied at equilibrium. Generation is then a forward pass from a clamped label;
inpainting and anomaly detection are settling with different clamps. Crucially, no code in the settling
or learning core changes: only the clamping pattern and the layer orientation differ.

\subsection{Masked and input-free settling}

The settling routine supports clamping the input fully (discriminative), leaving it free (generation),
or masking it so that some pixels are clamped and the rest are free (inpainting). When the input is
free, its update follows the same energy gradient, $-\phi'(s_0)\odot(W_0^{\!\top}\varepsilon_1)$, and a
binary mask keeps the known pixels fixed while relaxing the unknown ones. The same per-sample residual
energy $E = \tfrac12\sum_k\lVert\varepsilon_k\rVert^2$, evaluated at the predicted label, serves as an
anomaly score: an input the model cannot generate leaves high residual energy. This is the same
quantity that drives settling, repurposed as a detector.

\subsection{The output-dimension stability scaling}

A reusable stability finding emerged from the orientation flip and is worth recording as method. With a
$784$-dimensional output, the hidden-layer settling gradient (which contains the term
$\varepsilon_{\text{out}}\,W$) is roughly $\sqrt{784}\approx 28$ times larger than with a ten-dimensional
label output. A state learning rate that is well-behaved for the discriminative model therefore diverges
to NaN for the generative one. The state learning rate must be scaled down by roughly the square root of
the output dimension to remain stable. This is a general consequence of how the output error enters the
hidden-state gradient, not a tuning accident, and it generalises the choice of $\eta_x$ to other output
sizes.

\section{Optional Extensions: Incremental PC and EWC}

Two further methods are implemented and validated, not merely sketched, to position vanilla SO-PC
against close relatives.

\subsection{Incremental PC}

Incremental PC \citep{salvatori2024ipc} differs from standard SO-PC in \emph{when} the weights move.
Rather than settling fully and updating once at equilibrium, it updates the states \emph{and} the
weights after every inference step, from the same error vector. The update is simultaneous: the state
step uses the current weights and the weight step uses the pre-step states, both derived from one error
vector, so neither sees the other's change within a step. A clean correctness anchor pins this down ---
with the weight learning rate set to zero, one incremental step reduces \emph{exactly} to one standard
settling step, because incremental PC's state update is the SO update; it merely also moves the weights.

The consequential property is sensitivity to the weight learning rate. Incremental PC applies $T$ weight
updates per batch rather than one, so its effective weight learning rate is roughly $T$ times larger; at
the standard $\eta_w$ it diverges to NaN, and it needs a correspondingly smaller rate (on the order of
$\eta_w / T$) to be stable. The incremental path is restricted to the PyTorch backend by design: the
fused CUDA kernel settles \emph{without} touching the weights, so it cannot express an update that writes
the weights every step. The lesson --- incremental PC needs a $\sim T$-times smaller weight learning
rate, after which it is competitive or better --- is reported quantitatively in \Cref{sec:exp-depth}.

\subsection{Elastic Weight Consolidation}

To locate the absolute continual-learning numbers against a dedicated method, we add Elastic Weight
Consolidation \citep{kirkpatrick2017ewc} as a third class-incremental arm on the backprop side. After
each task we store the diagonal empirical Fisher information (the mean squared gradient of the task loss
at the converged parameters) and the optimal parameters, and later tasks carry the quadratic anchor
$\tfrac{\lambda}{2}\sum_i F_i(\theta_i - \theta_i^\star)^2$ that penalises moving parameters the earlier
task found important. EWC necessarily uses autograd (it is a backprop-side baseline, not PC) and slots
into the existing class-incremental harness as a method choice. A caveat carried over from the
fair-comparison work: in a regime too weak for the base learner to actually learn, EWC is correctly
inert, because there is nothing to forget. Its measured effect appears in \Cref{sec:exp-depth}.

\section{Reproducibility and Tooling}\label{sec:method-repro}

Reproducibility is treated as part of the method, not an afterthought.

\paragraph{Environment.} The project uses \texttt{uv} exclusively for environment management --- never
\texttt{pip} --- so the dependency set is locked and reconstructable with a single command. The CUDA
build is pinned to a toolkit-matched PyTorch (\texttt{cu126}) and requires nvcc with an MSVC toolchain
on Windows; the kernel is compiled from the real \texttt{.cu} source via JIT loading and cached. The
pure-PyTorch path requires no compilation at all, so the entire study --- minus the optional kernel ---
runs with only \texttt{torch} and \texttt{torchvision} installed.

\paragraph{Offline test suite.} A suite of forty tests validates the settling mathematics, the local
update, the energy formula, the generative and masked-settling paths, and the harness wiring. The tests
are hermetic and offline: they use small synthetic tensors with no MNIST download and run CPU-only, so
correctness can be re-verified anywhere without a GPU or network access. Earlier weak assertions (for
example, a training-loop test that only compared the last epoch to the first rather than checking the
trend) were tightened in the course of the project.

\paragraph{Autonomous hyperparameter search.} The Phase-4 search is a thin layer over
\texttt{train\_and\_eval}: grid and random search (with an optional Optuna backend) propose a config,
call the interface, and log the metrics, with no rewrite of the core. The single stable signature is
what makes the search ``almost free'' in engineering terms. As a sanity check, an autonomous search run
rediscovers the $\eta_x \leftrightarrow \eta_w$ stability interaction described in
\Cref{sec:protocol} --- a useful confirmation that the protocol's central confound is real and findable
without prior knowledge.

\paragraph{Figure and artefact pipeline.} Every quantitative claim is backed by a stored artefact:
benchmark JSON for the kernel timings and launch counts, per-seed result files for the alternating
sweep, and the figures of \Cref{ch:experiments}. Per-seed values are retained (not just aggregates) so
that the spread behind every interval can be inspected, which is essential given that several of the
study's headline findings are null and rest on the \emph{absence} of a separation rather than its
presence.
